% !TEX root = ../main.tex

\section{Preliminaries}
\paragraph{Notations.}
    We represent scalar discrete random variables over $N$ categories as one-hot column vectors, with ${\VC = \{\x \in \{0,1\}^N\colon \sum_{i=1}^N \x_i = 1\}}$ denoting the set of such vectors. Define $\text{Cat}(\cdot; \prior)$ as the categorical distribution over $N$ classes with probabilities given by $\prior \in \Delta$, where $\Delta$ denotes the $N$-simplex.
    % We further define the probability simplex $\Delta \vcentcolon= \{\z \in \RR^N\colon z \ge 0, \langle z, \vec{1} \rangle = 1 \}$ as the convex hull of $\VC$, where $\vec{1} \in \RR^N$ denotes the all-ones vector. 

\begin{tcolorbox}[colback=gray!20, colframe=gray!20, arc=2mm, boxrule=0pt, width=1\linewidth, boxsep=-1pt]
    \textbf{Generality of description.} We focus on Diffusion Language Models defined directly on discrete structures. While a fully general continuous-time formulation can be expressed in CTMC terms (SEDD; \citealp{lou2024discrete}), such generality is rarely used in practice: efficient training typically relies on specific forward processes. In particular, most works concentrate on two families: \textbf{absorbing (masked) noising} (MDLM; \citealp{sahoo2024simple}) and \textbf{uniform noising} (UDLM/Duo; \citealp{schiff2024discreteguidance, sahoo2025the}). These choices yield simpler sampling and practical training objectives. To keep the main text clear, we present only the absorbing and uniform processes and follow the notation of MDLM and Duo. For theoretical completeness, the general CTMC/SEDD formulation is considered in Appendix~\ref{app:details of considered processes}.
\end{tcolorbox}
    
    \subsection{Diffusion Language Models}
    \label{sec:discrete diffusion models}        
        We adopt the interpolating noise framework for diffusion language models introduced by~\citep{sahoo2024simple, sahoo2025the}, which considers a gradual interpolation between clean data samples $\x_0 \in \VC$, drawn from the data distribution $p^*(\x_0)$, and a simple reference distribution $\prior \in \Delta$. For a given clean input $\x_0 \in \VC$, the noising process is defined as follows:
        \begin{gather}
            p_{t|0}(\cdot|\x_0) = \text{Cat}(\cdot; \alpha_{t} \x_0 + (1-\alpha_{t})\prior),
        \end{gather}
        where $\alpha_{t} \in [0,1]$ is a strictly decreasing function in $t$ with $\alpha_0 \approx 1$ and $\alpha_1 \approx 0$. 
        % The evolution of marginals $p(\x_t)$ is described by a linear ordinary differential equation~\citep{anderson2012continuous, campbell2022continuous, lou2024discrete}:         
        % \begin{equation}
        %   \frac{dp_t}{dt} = Q_t p_t, \quad p_0 = p^*,
        % \label{eq:forward diffusion}
        % \end{equation}
        % where $Q_t \in \mathbb{R}^{N \times N}$ is the state transition matrix. 
        Two principal variants of the diffusion process have been proposed, namely the \emph{absorbing} and \emph{uniform} processes, which differ in their choice of terminal distribution $\prior$. We consider both of them in details below.
        
  \subsection{Absorbing diffusion (MDLM)}
\label{sec:mdlm}
Absorbing (masked) diffusion introduces a special mask token $\m \in \VC$ and set the terminal distribution to be concentrated on it, so $\prior = \m$. For $0 \le s < t \le 1$, the true posterior of absorbing process admits a simple two-case form $p(\x_s \mid \x_t, \x_0)=$:
\begin{equation}
\begin{cases}
\text{Cat}(\x_s; \x_t), & \x_t \neq \m,\\[2mm]
\text{Cat}\Big(\x_s;\, \dfrac{(1-\alpha_s)\m + (\alpha_s-\alpha_t)\x_0}{1-\alpha_t}\Big), & \x_t = \m.
\end{cases}
\end{equation}
MDLM parameterizes a denoiser $\widehat{x}_0(x_t,t) \in \Delta$ to approximate $\x_0$ and defines the reverse transition by substituting
$x_0 \leftarrow \widehat{x}_0(x_t,t)$ in the posterior:
\begin{gather}
  p_\theta(\x_s \mid \x_t)
  :=
  p\big(\x_s \mid \x_t, \x_0=\widehat{\x}_0(x_t,t)\big)
  =
  \\
  \begin{cases}
    \text{Cat}(\x_s; \x_t), & \x_t \neq \m,\\[2mm]
    \text{Cat}\Big(\x_s;\, \dfrac{(1-\alpha_s)\m + (\alpha_s-\alpha_t)\widehat{\x}_0(\x_t,t)}{1-\alpha_t}\Big), & \x_t = \m.
  \end{cases}
  \nonumber
\end{gather}
%
% In MDLM, this is implemented with the \textsc{subs} parameterization, which enforces that unmasked states are carried over and the denoiser never outputs the mask token, i.e. $\widehat{x}_0(x_t, t) = x_t$ if $x_t \neq m$. 
For a given denoiser $\widehat{\x}_0$, and starting distribution $\piv = \m$, this approximate posterior defines an iterative generative model. The denoiser $\widehat{\x}_0$ for his model is learned by optimizing negative evidence lower bound (NELBO), which in the continuous-time limit reduces to a weighted cross-entropy:
  \begin{Box1}{MDLM Loss, $\LC_{\text{MDLM}}(\widehat{x}_0, x_0) \vcentcolon=$}
    \vspace{-8pt}
    \begin{gather}
      \int_0^1 \, \EE_{p_{t|0}(\x_t \mid \x_0)} \left[
          \frac{\alpha_t'}{1-\alpha_t} \langle \log\widehat{\x}_0(\x_t, t), \x_0 \rangle
      \right] dt.
    \label{eq:mdlm loss}
    \end{gather}
  \end{Box1}
Note that since $\x_0$ is one-hot, $\log\langle \widehat{\x}_0(x_t,t), \x_0\rangle = \langle \log \widehat{\x}_0(\x_t,t), \x_0\rangle$.
  
  % For the absorbing process $\pi = m$ 
  
  % \begin{Box1}{MDLM Loss, $\LC_{\text{MDLM}}(\widehat{x}_0, x_0) \vcentcolon=$}
  %   \vspace{-8pt}
  %   \begin{gather}
  %     \int_0^1 \, \EE_{p_{t|0}(x_t|x_0)} \left[
  %         \lambda_{t} \langle \log\widehat{x}_0(x_t, t), x_0 \rangle
  %     \right]dt,
  %   \label{eq:mdlm loss}
  %   \end{gather}
  % \end{Box1}
  % %
  % where $\lambda_{t}$ is positive weighting function and $\log \widehat{x}_0(x_t, t)$ is element-wise logarithm of $\widehat{x}_0(x_t, t)$. Note that in the original work, the authors use the notation $\log \langle \widehat{x}_0(x_t, t), x_0 \rangle$, which is equivalent. Moreover, \citep{sahoo2024simple} show that this loss corresponds to the negative evidence lower bound (NELBO). As a result, at optimality, the learned model on whole data distribution $p^*(x_0)$: 
  % $$x_0^* = \arg\min_{\widehat{x}_0} \mathbb{E}_{p^*(x_0)}\left[\LC_{\text{MDLM}}(\widehat{x}_0, x_0)\right]$$
  % recovers the target reverse conditional density $p_{0\mid t}(x_0 \mid x_t).$
  % Moreover, \citet{sahoo2024simple} shows that the MDLM loss~\eqref{eq:mdlm loss} corresponds to the negative evidence lower bound (NELBO) in the limit as $T \to \infty$. Consequently, the optimal solution of this objective for the data distribution $p^*(x_0)$
  % \begin{gather}
  %   x_0^* = \arg\min_{\widehat{x}_0} \mathbb{E}_{p^*(x_0)}\LC_{\text{MDLM}}(\widehat{x}_0, x_0)
  %   \nonumber
  % \end{gather}
  % satisfies $x_0^*(x_t, t) = \EE_{p_{0 \mid t}(x_0 \mid x_t)} \left[x_0\right]= p_{0 \mid t}(x_0 \mid x_t)$.

\subsection{Uniform Process (Duo).}
  % Another key example is the \textit{uniform process}. 
  % where the terminal distribution is uniform, i.e., $\pi = \frac{1}{N}\vec{1}$. This is also achieved by choosing a specific base matrix $Q_{\text{uni}}$ (see Appendix~\ref{app:related works}), yielding the diffusion matrix $Q_t = -\frac{\alpha_t'}{N\alpha_t}(\vec{1}\vec{1}^{\top} - N I)$. 
  % A representative case of this process is the \textit{Uniform Diffusion Language Model} (UDLM; \citealp{schiff2024discreteguidance}), which adopts the same parameterization as MDLM. This leads to the simplified form of the SEDD loss~\eqref{eq:sedd loss}\footnote{For brevity, we omit the arguments and write $\widehat{x}_0 \vcentcolon= \widehat{x}_0(x_t, t)$.} for one data point $x_0 \sim p^*(x_0)$:
  % \begin{Box1}{UDLM Loss, $\LC_{\text{UDLM}}(\widehat{x}_0, x_0) \vcentcolon=$}
  %   \vspace{-8pt}
  %   \begin{gather}
  %   \label{eq:udlm loss}
  %     \int_0^1 \EE_{p_{t|0}(x_t|x_0)} \Bigg[
  %         \frac{\alpha_t'}{N\alpha_t} 
  %         \Bigg(\frac{1}{p_{t\mid0}(x_t \mid x_0)} - \frac{1}{p_{t \mid 0}(x_t \mid \widehat{x}_0)} \nonumber 
  %         \\ 
  %         \!\!\!\!\! - \! \sum_{y \neq x_t} \! \frac{p_{t\mid0}(y \!\mid \!x_0)}{p_{t\mid0}(x_t \! \mid \! x_0)}\log\frac{p_{t \mid 0}(x_t \!\mid\! \widehat{x}_0) p_{t \mid 0}(y \!\mid \!x_0)}{p_{t \mid 0}(y\!\mid\! \widehat{x}_0) p_{t \mid 0}(x_t \!\mid \! x_0)}\!\Bigg)\!\Bigg]dt.
  %   \end{gather}
  % \end{Box1}
  % %
  % \citet{sahoo2025the} further introduced the \textit{Diffusion Duality} (Duo) framework, connecting Gaussian and uniform processes. This leads to an extended parameterization $\widehat{x}_0\colon \Delta \times [0, T] \to \Delta$ and a minor modification of the UDLM loss~\eqref{eq:udlm loss}, while preserving its overall structure.

  % Additionally, \citet{schiff2024discreteguidance} show that the UDLM loss in equation~\eqref{eq:udlm loss} also corresponds to the NELBO in the limit as $T \to \infty$. Consequently, the optimal solution for the dadistributionion $p^*(x_0)$:
  % \begin{gather}
  %     x_0^* = \arg\min_{\widehat{x}_0} \mathbb{E}_{p^*(x_0)}\LC_{\text{UDLM}}(\widehat{x}_0, p^*) \nonumber
  % \end{gather} satisfies $x_0^*(x_t, t) =\EE_{p_{0 \mid t}(x_0 \mid x_t)} \left[x_0\right] = p_{0 \mid t}(x_0 \mid x_t)$.
  % (see Appendix~\ref{app:uniform process}
  The second important case is the \textit{uniform process}. Two notable models have been proposed in this setting: the \textit{Uniform Diffusion Language Model} (UDLM; \citealp{schiff2024discreteguidance}) and its extension, \textit{Diffusion Duality} (Duo; \citealp{sahoo2025the}). In this case $\piv = \vec{1}/N$. The forward corruption is given by:
  \begin{equation}
    p_{t|0}(\cdot \mid \x_0)
    = \text{Cat} \big(\cdot;\, \alpha_t \x_0 + (1-\alpha_t)\vec{1}/N\big),
  \end{equation}
  and corresponding posterior distribution is given by:
  \begin{gather}
    \label{eq:uniform true reverse posterior}    
    p(\x_s|\x_t, \x_0)
    = \text{Cat} \Bigg(\x_s;
        \frac{
            N\alpha_t \x_t \odot \x_0 + (\alpha_{t|s} - \alpha_t)\x_t
            }
            {
             N \alpha_t \langle \x_t, \x_0\rangle  + 1 - \alpha_t
            } \nonumber \\
        + \frac{
            (\alpha_s - \alpha_t)\x_0 + (1 - \alpha_{t|s})(1- \alpha_s)\vec{1} / N
            }
            {
             N \alpha_t \langle \x_t, \x_0\rangle  + 1 - \alpha_t
            }
        \Bigg),
\end{gather}
  where $\alpha_{t|s} = \alpha_t/\alpha_s$ and $\odot$ denotes element-wise product of vectors. Similar to MDLM, the UDLM parametrize denoiser $\widehat{x}_0(x_t, t)$ and consider approximate posterior:
  $$
    p_{\theta}(\x_s|\x_t, \x_0 = \widehat{x}(\x_t, t)).
  $$
  and derive the loss as NELBO as
  \begin{equation}
       \mathcal{L}_{\text{UDLM}}(\widehat{\x}_0,\x_0) \vcentcolon= \int_0^1 \EE_{p_{t|0}(\x_t|\x_0)} \left[g(\x_t,\x_0, \widehat{\x}_0(\x_t, t))\right] dt, \nonumber
  \end{equation}
  where the function $g(\x_t,\x_0, \widehat{x}_0(\x_t, t))$ is defined in \eqref{eq:udlm function}.

  Duo further introduces a theoretical connection between Gaussian and uniform diffusion processes. The authors consider distribution:
  \begin{equation}
    \tilde{p}_{t\mid 0}(w_t \mid \x_0) = \mathcal{N}(\tilde{\alpha}_t x_0, (1-\tilde{\alpha}_t^2)I), \nonumber
  \end{equation}
  where $\tilde{\alpha}_t$ denotes a rescaled diffusion schedule. And show that for $w_t \sim \tilde{p}_{t\mid 0}(w_t \mid \x_0)$:
  $$
    \x_t(w_t) \vcentcolon= \arg\max(w_t) \sim p_{t|0}^{\text{u}}(\x_t|\x_0)
  $$
  where $\arg\max(w_t)$ denotes a one-hot vector whose index corresponds to the maximum value in $w_t$.
  The corresponding loss becomes $\mathcal{L}_{\text{UDLM}}(\widehat{\x}_0,\x_0) \vcentcolon=$
  \begin{gather}
    \int_0^1 \EE_{\tilde{p}_{t\mid 0}(w_t\mid \x_0)} \left[g(\x_t(w_t),\x_0, \widehat{x}_0(\x_t(w_t), t))\right] dt, \nonumber
  \end{gather}
  However, the authors of Duo observe that this loss function exhibits high training variance, primarily due to the non-smooth nature of the $\argmax$ operation. To mitigate this issue, they propose reducing the variance by substituting the discrete model input $\x(w_t)$ with its continuous, soft approximation:
  \begin{equation}
      \x^{\tau}_{t}(w_t) = \mathrm{softmax}\left(\frac{w_t}{\tau}\right). \nonumber
  \end{equation} 
  % Note, that $\x_t(w_t) = \arg \max(\x_t^{\tau}) = \argmax(\softmax(\frac{w_t}{\tau}))$.
  This relaxation leads to an extended model parameterization $\widehat{x}_0 \colon \Delta \times [0, T] \to \Delta$. The resulting loss function for a data point $\x_0 \sim p^*(\x_0)$ is given by:
  \begin{Box1}{Duo Loss, $\LC_{\text{Duo}}(\widehat{x}_0, x_0) \vcentcolon=$}
    \vspace{-8pt}
    \begin{gather}
    \label{eq:duo loss}
      \int_0^1 \EE_{\tilde{p}_{t\mid 0}(w_t\mid x_0)} \left[g(\x_t(w_t),x_0, \widehat{x}_0(\x_t^{\tau}(w_t), t))\right]dt.
    \end{gather}
  \end{Box1}
  %
  Moreover, \citet{sahoo2025the} shows that, in the limit as $\tau \to 0^{+}$, this loss corresponds to the NELBO of UDLM. 
  
  % Consequently, at optimality and in this limit, the learned model on whole data distribution $p^*(x_0)$:
  % $$x_0^* = \arg\min_{\widehat{x}_0} \mathbb{E}_{p^*(x_0)}\LC_{\text{Duo}}(\widehat{x}_0, x_0)$$
  % recovers the target reverse conditional distribution $p_{0\mid t}(x_0 \mid x_t).$
  

\paragraph{Sequence-Level Discrete Diffusion.} Extension for both absorbing and uniform diffusion on a sequence $\x^{1:L}$ of $L$ tokens is made by considering forward and posterior distribution factorized independently. Specifically, the forward process is given by:
\begin{gather}
    p_{t|0}(\x_t^{1:L} \mid \x_0^{1:L})
    = \prod_{l=1}^L p_{t|0}(\x_t^l \mid \x_0^l).
    \nonumber
\end{gather}
The approximate posterior given by:
\begin{gather}
    p^{\theta}_{s|t}(\x_s^{1:L}|\x_t^{1:L}) = \prod_{l}p_{s|t}(\x_s^l|\x_t^l,\widehat{x}_0(\x_t^{1:L}, t)),
    \nonumber
\end{gather}
where denoiser is $\widehat{x}_0: \VC^{L} \times [0, 1] \rightarrow \Delta^{L}$ and all considered objectives decomposes on sum over tokens.
% $$
%     \mathcal{L}(\widehat{\x}_0, \x_0^{1:L}) = \sum_{l=1}^L \mathcal{L}(\widehat{\x}_0, \x_0^l)
% $$
  % In practical applications, we are interested in generating sequences of some length. A direct extension of the diffusion process to this setting would require modeling over an exponentially large space, resulting in a diffusion matrix of intractable size. To overcome this limitation, following~\citep{lou2024discrete}, we consider diffusion processes defined by sparse diffusion matrices that perturb tokens independently across sequence positions. Under this assumption, the forward diffusion matrix factorizes across tokens, yielding both a factorized conditional distribution and, crucially, a loss function that decomposes as a sum of token-wise losses. As a result, the per-token objectives in equations~\eqref{eq:sedd loss}, \eqref{eq:mdlm loss}, and~\eqref{eq:duo loss} can be directly applied to sequence modeling.
% \begin{gather}
%     p_{t \mid 0}(x_t^{1:L} \mid x_0^{1:L}) = \prod_{i=1}^L p_{t \mid 0}(x_t^{i}\mid x_0^{i}), \nonumber
% \end{gather}
% This structure induces a corresponding model parameterization
% $\widehat{s}\colon \XC^L \times [0, T] \to \RR_{+}^{N \times L}$,
% and yields the following SEDD loss:
% \begin{gather}
%     \textit{Sequence SEDD loss}: \label{eq:sequence sedd loss}\\
%     \LC_{\text{SEDD}}(\widehat{s}, p) \vcentcolon=
%     T \EE_{x_0^{1:L}, x_t^{1:L}, t}
%     \sum_{i=1}^L \sum_{y \neq x_t^i}
%     Q_t(x_t^i, y)\, g_{\text{SEDD}}(\widehat{s}), \nonumber \\
%     x_0^{1:L} \sim p(x_0^{1:L}), \quad
%     x_t^{1:L} \sim p_{t \mid 0}(x_t^{1:L} \mid x_0^{1:L}), \quad
%     t \sim \UC([0, T]). \nonumber
% \end{gather}
% The same factorization extends naturally to the (M/U)DLM setting, resulting in the model parameterization $\widehat{x}_0\colon \XC^L \times [0, T] \to \Delta^L$. Analogously, for the Duo model the parameterization takes the form $\widehat{x}_0\colon \Delta^L \times [0, T] \to \Delta^L$. Under this formulation, the loss function is given by:
% \begin{gather}
%     \textit{Sequence (M/U)DLM loss}: \label{eq:sequence (m/u)dlm loss}\\
%     \LC_{\text{(M/U)DLM}}(\widehat{x}_0, p) \vcentcolon=
%     T \EE_{x_0^{1:L}, x_t^{1:L}, t}
%     \sum_{i=1}^L \sum_{y \neq x_t^i}
%     Q_t(x_t^i, y)\, g_{\text{(M/U)DLM}}(\widehat{x}_0), \nonumber \\
%     x_0^{1:L} \sim p(x_0^{1:L}), \quad
%     x_t^{1:L} \sim p_{t \mid 0}(x_t^{1:L} \mid x_0^{1:L}), \quad
%     t \sim \UC([0, T]). \nonumber
% \end{gather}

\subsection{Distillation of Diffusion Models} 

\paragraph{Acceleration of diffusion models via distillation.}
  Diffusion language models share the same problem of long iterative inference as continuous state space diffusion models. One of the primary approaches to mitigating this issue in continuous state space models is the use of a procedure known as \emph{distillation}. This technique involves training a new model, referred to as the \textbf{student}, to approximate the sampling behavior of the original diffusion model, termed the \textbf{teacher}, but with significantly fewer inference steps. For continuous space diffusion models, a huge number of different distillation techniques were developed~\citep{luhman2021knowledge, salimans2022progressive, song2023consistency}. 
  % In this work, we focus on a particular class of such methods, known as \emph{Inverse Distillation}.

\paragraph{Inverse Distillation of Diffusion Models.} 
  % Although discrete diffusion models have demonstrated strong performance in text generation tasks, their reliance on multi-step inference substantially limits their practical applicability. In contrast, a large body of work has focused on accelerating diffusion models in continuous spaces, among which distillation-based approaches have emerged as particularly effective~\citep{luhman2021knowledge,salimans2022progressive,song2023consistency}. These approaches aim to distill a pre-trained multi-step diffusion model into a few-step student generator capable of producing high-quality samples. A notable instantiation of this approach involves training an auxiliary or “fake” model as part of the distillation process to support the training of the few-step student generator. While these methods were initially developed for different classes of models, they can be unified under the broader framework of \textit{inverse distillation}. This approach was first introduced in the Inverse Bridge Matching Distillation framework~\citep{gushchin2025inverse}, where the authors derived a distillation objective based on an inverse Kullback–Leibler (KL) divergence. Subsequently, \citet{kornilov2025universal} demonstrated that several prior works~\citep{zhou2024score, zhou2024adversarial, huang2024flow} also implicitly can be described within this framework. Nevertheless, these connections have thus far been explored primarily in the context of continuous-space diffusion models and presented as auxiliary theoretical insights rather than as core methodological principles. In contrast, our work explicitly adopts the inverse distillation perspective and extends it to the discrete diffusion setting, positioning it as a central component of our approach.
  A particularly prominent class of distillation techniques involves the integration of an auxiliary diffusion model, referred to as the \textbf{fake} model, into the training process~\citep{yin2024one,yin2024improved,zhou2024score, zhou2024adversarial, huang2024flow, gushchin2025inverse, kornilov2025universal}. One of the branches of this work is the direction of 
  % In this work, we build upon the concept of 
  \textit{Inverse Distillation}, originally introduced in the Inverse Bridge Matching Distillation framework~\citep{gushchin2025inverse} for a special type of diffusion models called ``Diffusion Bridge Models'' and further developed by~\citet{kornilov2025universal} for the general case of \underline{continuous space} flow and diffusion models. 

\paragraph{Teacher as an optimizer.}
  Let $p^*(x_0)$ denote the data distribution and let $\LC_{\text{cont.}}(f; x_0)$ be the per-example training loss
  used to fit a continuous diffusion model $f$ from data.
  The pretrained \emph{teacher} model is defined as the optimizer of this objective:
  \begin{equation}
    f^* = \arg\min_f \ \mathbb{E}_{p^*(x_0)}\big[ \LC_{\text{cont.}}(f; x_0)\big]. \nonumber
    \label{eq:teacher_objective_cont}
  \end{equation}

% One important instance of \eqref{eq:teacher_objective_cont} is \emph{diffusion/flow matching}
% \citep{kornilov2025}, where one matches a model field $\hat f(x_t,t)$ to the conditional field induced by
% the forward process started at $x_0$:
% \begin{equation}
%   L_d(f; x_0)
%   := \mathbb{E}_{t\sim \mathrm{Unif}[0,T]}\,
%      \mathbb{E}_{x_t\sim p_{t|0}(\cdot|x_0)}
%      \big\|\hat f(x_t,t) - f_{p_0}(x_t,t\,|\,x_0)\big\|^2 .
%   \label{eq:matching_loss}
% \end{equation}
% In this case one can take $L_{\text{teach}} \equiv L_d$.

\paragraph{Inverse objective.}
Inverse distillation optimizes a \emph{student distribution} $p_\theta(x_0)$ such that the fixed teacher $f^*$ remains optimal under $p_\theta$. Specifically, inverse distillation considers objective $\LC_{\text{inv}}(\theta) :=$
\begin{gather}
     \!\!\mathbb{E}_{p_{\theta}(x_0)} \!\left[\LC_{\text{cont.}}(f^*\!\!, x_0)\right]  - \min_{\widehat{f}} \mathbb{E}_{p_{\theta}(x_0)} \! \left[\LC_{\text{cont.}}(\widehat{f}, x_0)\right], 
     \label{eq:inverse optimization problem} 
\end{gather}

  % \begin{Box1}{Inverse Distillation, $\LC_{\text{inv}}(\theta) \vcentcolon=$}
  %   \begin{equation}\label{eq:inverse optimization problem} 
  %       \!\!\! \mathbb{E}_{p_{\theta}(x_0)} \!\left[\LC_{\text{d/f}}(f^*\!\!, x_0)\right]  - \min_{\widehat{f}} \mathbb{E}_{p_{\theta}(x_0)} \! \left[\LC_{\text{d/f}}(\widehat{f}, x_0)\right]\!, \!\!
  %   \end{equation}
  % \end{Box1}
  % \begin{equation}
  %   \LC_{\text{inv}}(\theta)
  %   :=
  %   \mathbb{E}_{x_0\sim p_\theta}\big[L_{\text{teach}}(f^\*; x_0)\big]
  %   \;-\;
  %   \min_f \ \mathbb{E}_{x_0\sim p_\theta}\big[L_{\text{teach}}(f; x_0)\big],
  %   \label{eq:inverse_distillation_cont}
  % \end{equation}
  where the inner minimizer is the \emph{fake} model trained on samples from $p_\theta$. 
  By construction, the inverse loss satisfies $\LC_{\text{inv}}(\theta) \geq 0$ and for $p_\theta(x_0) = p^*(x_0)$ the minimum value $0$ is attained. However, this formulation does not, in general, guarantee the uniqueness of the minimizer. Consequently, the optimization procedure may converge to suboptimal solutions.
  % By definition, $\LC_{\text{inv}}(\theta)\ge 0$, and for $p_\theta(x_0) = p^*(x_0)$ the minimum value $0$ is attained.

  % Since \rewriteit{there is no PF-ODE analog for discrete diffusion}, 
  % We focus our attention on the distillation using an auxiliary \textbf{fake} model. One of the branches of this work is the direction of 
  % % In this work, we build upon the concept of 
  % \textit{Inverse Distillation}, originally introduced in the Inverse Bridge Matching Distillation framework~\citep{gushchin2025inverse} for a special type of diffusion models called "Diffusion Bridge Models" and further developed by~\citet{kornilov2025universal} for the general case of \underline{continuous space} flow and diffusion models. 
  
  % However, prior studies have primarily explored this idea in the context of continuous diffusion models, while our work explicitly extends the inverse distillation to the discrete setting.
  % This paradigm considers a unified learning objective for flows and diffusions based on the unification of conditional score and conditional flow matchings. Consider a \underline{continuous space} (specifically D-dimensional real space $\mathbb{R}^D$) probability path $\{p_t \}_{t\in 0, T}$ on the time interval $[0, T]$ that transforms data distribution $p^*_0$ into noise distribution $p_T$. The path itself is built as a mixture of simple conditional paths $\{p_t(\cdot|x_t) \}_{t\in [0, T]}$ conditioned on samples $x_0 \sim p(x_0)$:
  % \begin{gather}
  %   p_t(x_t) = \int_{\mathbb{R}^D} p_t(x_t|x_0)p(x_0)dx_0.
  %   \nonumber
  % \end{gather}
  % In turn, with each conditional path $p_t(x_t|x_0)$ there is associated a function $f^{p^0}(\cdot|x_0)$ which recovers it (i.e., score, or ODE/SDE drift functions of diffusion started from data point $x_0$). To fit a function $f^{p_0}$ which recovers $p^*(x_0)$ from the noise $p_T$ the authors propose a unified matching loss:
  %   \begin{Box1}{\rewriteit{Diffusion/Flow Matching}, $\LC_{\text{d}} \vcentcolon=$}
  %   \begin{equation}\label{eq:inverse optimization problem} 
  %       \mathbb{E}_{p(x_t| x_0)} \|\widehat{f}(x_t) - \widehat{f}_t^{p_0}(x_t|x_0)\|^2
  %   \end{equation}
  % \end{Box1}
  % The model $f^*$  trained by the minimization of this loss is called \textbf{teacher} model:
  % \begin{gather}
  %     f^*~=~\argmin_{\widehat{f}} \mathbb{E}_{p^*(x_0)} \left[\LC_{\text{d}}(\widehat{f}, x_0)\right].
  %     \nonumber
  % \end{gather}
  % This paradigm considers a parametric family of optimization problems defined by an objective function $\LC(\widehat{f}, \theta)$, where $\widehat{f}$ denotes the optimization variable and $\theta$ represents the problem parameters. We define the \underline{forward optimization problem} as finding the optimal solution $f^*$ for given parameters $\theta$, that is, $f^*~=~\argmin_{\widehat{f}} \LC(\widehat{f}, \theta)$. 
  % Then the authors propose a general way of distilling this model into a one-step generator $G_{\theta}$, which parametrizes the distribution $p_{\theta}(x_0)$ using the following objective:
  % Conversely, the inverse optimization problem assumes access to a desired solution $f^*$ and aims to recover parameters $\theta^*$ such that $f^*$ solves the corresponding forward problem $f^* = \arg\min_{\widehat{f}} \LC(\widehat{f}, \theta^*)$. A standard approach to solving the inverse problem is to minimize the following objective~\citep{gushchin2025inverse, kornilov2025universal}:
  % where model $\widehat{f}$ is called \textbf{fake} model. By construction, the inverse loss satisfies $\LC_{\text{inv}}(\theta) \geq 0$. Furthermore, for the case of $p_{\theta}(x_0) = p^*(x_0)$ (i.e., $G_{\theta}$ produces a data distribution in one step), the loss attains its minimum value of 0. 

  % \textcolor{red}{Write gradient formula for inverse loss and say that everything is easy in case of continious diffusion. Or do it in the method section?}

  While such a general description of the inverse distillation was proposed in~\citep{kornilov2025universal} the particular cases of this distillation were introduced before, specifically for continuous diffusion (Score identity Distillation, SiD) in~\citep{zhou2024score} and specifically for continuous flows (Flow Generator Matching, FGM) in~\citep{huang2024flow}. 

  \paragraph{Connection to this work.}
    In \emph{continuous} space, the model distribution $p_\theta$ is typically reparameterized, allowing the optimization of $\LC_{\text{inv}}(\theta)$ in Equation~\eqref{eq:inverse optimization problem} via backpropagation through sampled trajectories. In contrast, extending this approach to the \emph{discrete} domain presents several challenges, stemming from the inherent characteristics of discrete spaces. These challenges include (i) \underline{theoretical issues} concerning the validity of the training objective, particularly due to the potential non-uniqueness of its minimizers, and (ii) \underline{practical difficulties} associated with the non-smooth nature of discrete space. Our objective is to extend the inverse distillation framework to the widely studied class of \emph{Diffusion Language Models} by addressing both of these challenges, as discussed in~(\wasyparagraph\ref{sec:idlm}).
  % While the loss attains its minimum value of $0$ at the trivial solution $\theta = \theta^*$, this formulation does not guarantee the uniqueness of the solution. As a result, the optimization process may converge to suboptimal solutions.
  
  % In this work, we adopt the \textit{Inverse Optimization} framework~\citep{chan2025inverse} as the foundation for our main method. This paradigm considers a parametric family of optimization problems defined by an objective function $\LC(\widehat{f}, \theta)$, where $\widehat{f}$ denotes the optimization variable and $\theta$ represents the problem parameters. The forward problem seeks the optimal solution $f^*$ for a given set of parameters $\theta$, i.e., $f^* = \arg\min_{\widehat{f}} \LC(\widehat{f}, \theta)$. Conversely, the inverse optimization problem assumes access to a desired solution $f^*$ and aims to recover parameters $\theta^*$ such that $f^*$ solves the corresponding forward problem $f^* = \arg\min_{\widehat{f}} \LC(\widehat{f}, \theta^*)$. A standard approach to solving the inverse problem is to minimize the following objective:
  % \begin{Box1}{Inverse Optimization Objective}
  %   \begin{equation}
  %     \label{eq:inverse optimization problem} 
  %     \LC_{\text{inv}}(\theta) \vcentcolon= \LC(f^*, \theta) - \min\nolimits_{\widehat{f}} \LC(\widehat{f}, \theta).
  %   \end{equation}
  % \end{Box1}
  % %
  % By construction, the inverse loss satisfies $\LC_{\text{inv}}(\theta) \geq 0$. While the loss attains its minimum value of $0$ at the trivial solution $\theta = \theta^*$, this formulation does not guarantee the uniqueness of the solution. As a result, the optimization process may converge to suboptimal solutions.

  % It is worth noting that \citet{kornilov2025universal} establish a connection between their proposed loss function and an inverse optimization formulation. However, this connection is developed in the context of continuous-space diffusion models and is presented primarily as an auxiliary theoretical insight rather than a central methodological component. In contrast, our work builds directly upon this inverse optimization perspective and extends it to the setting of discrete diffusion models.
